% Generated draft: requires mathematical and source review.
\documentclass[11pt]{article}
\usepackage[margin=25mm]{geometry}
\usepackage{fontspec}
\setmainfont{Latin Modern Roman}
\usepackage{amsmath,amssymb,mathtools,unicode-math}
\setmathfont{Latin Modern Math}
\usepackage{enumitem,needspace,xurl,hyperref}
\setlength{\parindent}{0pt}
\setlength{\parskip}{5pt}
\setlength{\emergencystretch}{3em}
\newcommand{\N}{\mathbb{N}}
\newcommand{\Z}{\mathbb{Z}}
\newcommand{\Q}{\mathbb{Q}}
\newcommand{\R}{\mathbb{R}}
\newcommand{\C}{\mathbb{C}}
\newcommand{\F}{\mathbb{F}}
\newcommand{\A}{\mathbb{A}}
\newcommand{\PP}{\mathbb{P}}
\newcommand{\E}{\mathbb{E}}
\newcommand{\eps}{\varepsilon}
\newcommand{\dd}{\,\mathrm{d}}
\begin{document}
\section*{UnsolvedMath ID: 11000082}
\textbf{Source problem code:} catalogue entry
\subsection{Definitions and mathematical statement}
Question 7.6 asks whether $\operatorname{spec}(\mathcal I_g(k))$ has bounded multiplicity for $k\geq3$. Neither the subgroup notation $\mathcal I_g(k)$ nor the spectrum and its counting convention are defined in the supplied excerpt.

In the natural Teichm\"uller-length interpretation, a pseudo-Anosov $\varphi$ with stretch factor $\lambda(\varphi)>1$ has length $\ell(\varphi)=\log\lambda(\varphi)$. After specifying which conjugacy classes or closed geodesics are counted, write $m_{g,k}(t)$ for the number of those classes from $\mathcal I_g(k)$ with length $t$. Bounded multiplicity for each fixed $k$ would mean
\[
 \forall k\geq3\quad\exists C_{g,k}<\infty\quad
 \forall t>0,\qquad m_{g,k}(t)\leq C_{g,k}.
\]
The source does not prescribe a bound uniform in $g$ or $k$. This formulation is conditional on the missing spectrum definition. Counting conjugacy inside $\mathcal I_g(k)$ rather than inside $\operatorname{Mod}_g$, distinguishing inverse classes, and counting primitive classes rather than all powers can change multiplicity. A set of length values alone does not retain multiplicity.

The source suggests considering instead the simple length spectrum: the part of $\operatorname{spec}(\operatorname{Mod}_g)$ arising from pseudo-Anosovs represented by non-self-intersecting geodesic loops in $\mathcal M_g$. Simplicity here concerns loops in moduli space, not simple closed curves on the underlying surface. The excerpt does not state a bounded-multiplicity theorem for this simple spectrum or define the precise geodesic/orbifold convention.
\subsection{Short English statement}
For each filtration level $k\geq3$, is there a uniform bound on how many counted pseudo-Anosov classes have any one spectral length? The missing subgroup and multiplicity conventions are essential.
\subsection{Sources}
\begin{itemize}
\item Local registry: \url{data/unsolvedmath/problems.json}, permanent UnsolvedMath ID $11000082$, exact source code \texttt{catalogue entry}.
\item Farb, \emph{Problems on Mapping Class Groups and Related Topics}, Question 7.6, PDF page 58: \url{https://www.math.uchicago.edu/~farb/papers/mcgbook.pdf}. Supplied access date: 2026-07-29; extraction: pdf-font-structure; page-disambiguated numbering among 329 headings.
\item No additional reference is supplied. The registry explicitly says it cannot verify the exact invariant; preceding definitions are absent.
\end{itemize}
\end{document}
